% ============================================================
% Section 52: 德拜模型——声子总数的温度依赖与德拜温度计算
% ============================================================
\newpage
\section{固体理论：德拜模型——声子总数的温度依赖与德拜温度计算}
\label{sec:debye-phonon-number}

\begin{modelbox}[题目：德拜模型下声子总数的温度依赖与德拜温度]
在三维晶体中，利用德拜模型：
\begin{enumerate}
    \item 证明高温时，$0\sim\omega_D$ 范围内声子总数与温度 $T$ 成正比；
    \item 证明低温下，$0\sim\omega_D$ 范围内声子总数与温度 $T^3$ 成正比；
    \item 若已知一面心立方单原子晶体的晶格常数 $a=5\,\text{\AA}$，横向弹性波速 $v_t=3.54\times10^5\,\mathrm{cm/s}$，纵向弹性波速度 $v_l=4.92\times10^5\,\mathrm{cm/s}$，试计算其德拜温度。
\end{enumerate}
\end{modelbox}

\begin{tipbox}[考点快速分析]
\begin{itemize}
    \item \textbf{德拜模型}：把晶体视为各向同性连续弹性介质，色散线性 $\omega=vq$，存在截止频率 $\omega_D$
    \item \textbf{模式密度}：$D(\omega)\propto\omega^2$（三维），这是低温 $T^3$ 律的根源
    \item \textbf{普朗克分布}：$n(\omega)=(e^{\hbar\omega/k_BT}-1)^{-1}$
    \item \textbf{高温极限}：$k_BT\gg\hbar\omega$，$n(\omega)\approx k_BT/\hbar\omega$，每个模式声子数与 $T$ 成正比
    \item \textbf{低温极限}：$k_BT\ll\hbar\omega_D$，只有 $\omega\lesssim k_BT/\hbar$ 的模式被激发，有效模式数 $\propto T^3$
    \item \textbf{德拜温度}：$\Theta_D=\hbar\omega_D/k_B$，由总模式数 $3N$ 的截断条件确定
\end{itemize}
\end{tipbox}

\begin{derivationbox}[标准解题过程]
\textbf{Step 0: 德拜模型的模式密度}

三维晶体中有一支纵波和两支横波，色散均为线性 $\omega=v_iq$。在波矢空间，模式密度为 $V/(2\pi)^3$。对每一支波，频率 $\omega\sim\omega+d\omega$ 内的模式数为
\begin{equation}
D_i(\omega)\,d\omega = \frac{V}{(2\pi)^3}\cdot4\pi q^2\,dq = \frac{V}{2\pi^2}\frac{\omega^2}{v_i^3}\,d\omega.
\end{equation}

三支波的总模式密度：
\begin{equation}
D(\omega) = \frac{V}{2\pi^2}\Bigl(\frac{1}{v_l^3}+\frac{2}{v_t^3}\Bigr)\omega^2 \equiv \frac{3V}{2\pi^2v_p^3}\,\omega^2,
\label{eq:debye-Domega}
\end{equation}
其中定义了平均声速 $v_p$ 满足
\begin{equation}
\frac{3}{v_p^3}=\frac{1}{v_l^3}+\frac{2}{v_t^3}.
\end{equation}

频率为 $\omega$ 的平均声子数由普朗克分布给出：
\begin{equation}
n(\omega)=\frac{1}{e^{\hbar\omega/k_BT}-1}.
\label{eq:planck}
\end{equation}

\textit{关于 $n(\omega)$ 与 $D(\omega)$ 的物理角色}：
\begin{itemize}
    \item $D(\omega)$ 是\textbf{系统的固有属性}（单位频率内的模式数），由晶体结构、维度、色散关系决定，与温度无关。类比：宿舍楼的``房间容量表''——每层楼有多少个房间。
    \item $n(\omega)$ 是\textbf{热平衡下的统计占据}（每个模式上的平均声子数），由温度和玻色子统计决定。类比：每间宿舍平均住了多少学生。
    \item 二者相乘再积分：$\displaystyle N_{\text{ph}}=\int n(\omega)D(\omega)\,d\omega$，即``每个模式的人数 × 模式总数''，得到整个系统的总声子数。类似地，$U=\int\hbar\omega\,n(\omega)D(\omega)\,d\omega$ 只是多了一个每个声子的能量因子 $\hbar\omega$。
\end{itemize}

总声子数：
\begin{equation}
N_{\text{ph}}=\int_0^{\omega_D}n(\omega)D(\omega)\,d\omega
=\frac{3V}{2\pi^2v_p^3}\int_0^{\omega_D}\frac{\omega^2}{e^{\hbar\omega/k_BT}-1}\,d\omega.
\label{eq:Nph-integral}
\end{equation}

令 $x=\hbar\omega/k_BT$，则 $\omega=k_BTx/\hbar$，$d\omega=k_BT\,dx/\hbar$，上限 $x_D=\hbar\omega_D/k_BT=\Theta_D/T$：
\begin{equation}
N_{\text{ph}}=\frac{3V}{2\pi^2v_p^3}\Bigl(\frac{k_BT}{\hbar}\Bigr)^3\int_0^{\Theta_D/T}\frac{x^2}{e^x-1}\,dx.
\label{eq:Nph-dimensionless}
\end{equation}

\textbf{Step 1: 高温极限 ($T\gg\Theta_D$)}

此时 $x_D=\Theta_D/T\to0$，积分上限很小。在 $x\ll1$ 时 $e^x-1\approx x$，故
\begin{equation}
\int_0^{\Theta_D/T}\frac{x^2}{e^x-1}\,dx \approx \int_0^{\Theta_D/T}x\,dx = \frac{1}{2}\Bigl(\frac{\Theta_D}{T}\Bigr)^2.
\end{equation}

代入 \eqref{eq:Nph-dimensionless}：
\begin{equation}
N_{\text{ph}} \approx \frac{3V}{2\pi^2v_p^3}\Bigl(\frac{k_BT}{\hbar}\Bigr)^3\cdot\frac{1}{2}\Bigl(\frac{\Theta_D}{T}\Bigr)^2
=\frac{3Vk_B^3\Theta_D^2}{4\pi^2\hbar^3v_p^3}\cdot T.
\end{equation}

系数为常数，故 \textbf{$N_{\text{ph}}\propto T$}（高温）。

\textit{物理图像}：高温下 $k_BT\gg\hbar\omega_D$，所有 $3N$ 个模式上的声子数都很大，$n(\omega)\approx k_BT/\hbar\omega\gg1$。总声子数正比于被激发的模式数（常数 $3N$）乘以每个模式的平均声子数（$\propto T$）。

\textbf{Step 2: 低温极限 ($T\ll\Theta_D$)}

此时 $x_D=\Theta_D/T\to\infty$，积分上限可延拓至 $\infty$。利用标准积分
\begin{equation}
\int_0^{\infty}\frac{x^2}{e^x-1}\,dx = \Gamma(3)\zeta(3)=2\zeta(3)\approx2.404,
\end{equation}
其中 $\Gamma(3)=2!=2$，$\zeta(3)\approx1.202$（Ap\'ery常数）。

代入 \eqref{eq:Nph-dimensionless}：
\begin{equation}
N_{\text{ph}} = \frac{3V}{2\pi^2v_p^3}\Bigl(\frac{k_BT}{\hbar}\Bigr)^3\cdot2\zeta(3)
=\frac{3V\zeta(3)k_B^3}{\pi^2\hbar^3v_p^3}\cdot T^3.
\end{equation}

系数为常数，故 \textbf{$N_{\text{ph}}\propto T^3$}（低温）。

\textit{物理图像——三步拆解}：

\textbf{第一步：哪些模式被激发？}

看普朗克分布 $n(\omega)=1/(e^{\hbar\omega/k_BT}-1)$。当 $\hbar\omega\gg k_BT$ 时，指数项 $e^{\hbar\omega/k_BT}\gg1$，故 $n(\omega)\approx e^{-\hbar\omega/k_BT}\to0$——这些高频模式上几乎没有声子，相当于``冻结''了。只有当 $\hbar\omega\lesssim k_BT$ 时，$n(\omega)\sim O(1)$，模式才被显著激发。

因此，低温下有效的频率范围被热能量 $k_BT$ 截断在
\begin{equation}
\omega_{\text{cutoff}}\sim\frac{k_BT}{\hbar}.
\end{equation}

\textbf{第二步：这个频率范围内有多少个模式？}

模式密度 $D(\omega)\propto\omega^2$。被激发的模式总数 = 对 $D(\omega)$ 从 $0$ 到 $\omega_{\text{cutoff}}$ 积分：
\begin{equation}
N_{\text{modes}}\sim\int_0^{k_BT/\hbar}\omega^2\,d\omega\sim\Bigl(\frac{k_BT}{\hbar}\Bigr)^3\propto T^3.
\end{equation}

这是 $T^3$ 的第一个来源——\textbf{``有资格被激发的房间数''随温度三次方增长}。

\textbf{第三步：每个被激发的模式上有多少个声子？}

在截断频率 $\omega\sim k_BT/\hbar$ 处：
\begin{equation}
n\sim\frac{1}{e^{1}-1}\approx0.58\sim O(1).
\end{equation}

在更低的频率处 $n$ 更大，但大部分被积模式的 $n$ 在 $0.5\sim2$ 之间，平均下来是一个与温度无关的常数（约为 $1.2$，由积分中的 $2\zeta(3)/\text{积分上限}^3$ 决定）。

因此：
\begin{equation}
\underbrace{N_{\text{ph}}}_{\propto T^3}=\underbrace{N_{\text{modes}}}_{\propto T^3}\times\underbrace{\langle n\rangle}_{\sim O(1)}.
\end{equation}

$\boxed{T^3}$ 的本质：低温下热能量 $k_BT$ 很小，只能``唤醒''频率低于 $k_BT/\hbar$ 的模式。而三维空间中模式密度 $D(\omega)\propto\omega^2$，低频模式极其稀少——就像一个倒立的锥体，越靠近底部（低频）空间越小。温度升高时，$k_BT$ 这个``唤醒阈值''上移，允许被激发的模式数按体积（$\propto T^3$）增长。

\textbf{Step 3: 计算面心立方晶体的德拜温度}

德拜频率 $\omega_D$ 由总模式数等于晶体总自由度 $3N$ 确定：
\begin{equation}
\int_0^{\omega_D}D(\omega)\,d\omega = 3N.
\end{equation}

代入 \eqref{eq:debye-Domega}：
\begin{equation}
\frac{3V}{2\pi^2v_p^3}\int_0^{\omega_D}\omega^2\,d\omega = \frac{V\omega_D^3}{2\pi^2v_p^3}=3N
\quad\Longrightarrow\quad
\omega_D=\Bigl(\frac{6\pi^2N}{V}\Bigr)^{1/3}v_p.
\end{equation}

\textit{平均声速 $v_p$ 的计算}：
\begin{align}
\frac{3}{v_p^3}&=\frac{1}{v_l^3}+\frac{2}{v_t^3}
=\frac{1}{(4.92\times10^5)^3}+\frac{2}{(3.54\times10^5)^3}\\
&\approx 8.40\times10^{-18}+2\times2.25\times10^{-17}
=5.34\times10^{-17}\;(\mathrm{s^3/cm^3}),\\[4pt]
v_p&=\Bigl(\frac{3}{5.34\times10^{-17}}\Bigr)^{1/3}\approx3.83\times10^5\;\mathrm{cm/s}.
\end{align}

\textit{面心立方的原子数密度}：晶胞含 4 个原子，故
\begin{equation}
\frac{N}{V}=\frac{4}{a^3}=\frac{4}{(5\times10^{-8}\,\mathrm{cm})^3}=\frac{4}{1.25\times10^{-22}}=3.20\times10^{22}\;\mathrm{cm^{-3}}.
\end{equation}

\textit{德拜频率}：
\begin{align}
\omega_D&=\bigl(6\pi^2\cdot3.20\times10^{22}\bigr)^{1/3}\times3.83\times10^5\\
&\approx(1.895\times10^{24})^{1/3}\times3.83\times10^5\\
&\approx1.24\times10^8\times3.83\times10^5\approx4.75\times10^{13}\;\mathrm{rad/s}.
\end{align}

\textit{德拜温度}：
\begin{equation}
\Theta_D=\frac{\hbar\omega_D}{k_B}
=\frac{1.055\times10^{-34}\times4.75\times10^{13}}{1.38\times10^{-23}}
\approx\boxed{363\,\mathrm{K}}.
\end{equation}
（与参考值 $362\,\mathrm{K}$ 一致，中间舍入差异。）
\end{derivationbox}

\begin{formulabox}[答案汇总]
\begin{center}
\begin{tabular}{cl}
\toprule
问题 & 结论 \\
\midrule
(1) 高温声子数 & $N_{\text{ph}}\propto T$ \\
(2) 低温声子数 & $N_{\text{ph}}\propto T^3$ \\
(3) 德拜温度 & $\Theta_D\approx362\,\mathrm{K}$（FCC，$a=5\,\text{\AA}$，$v_l=4.92\times10^5\,\mathrm{cm/s}$，$v_t=3.54\times10^5\,\mathrm{cm/s}$） \\
\bottomrule
\end{tabular}
\end{center}
\end{formulabox}

\begin{insightbox}[物理图像与概念深化]
\textbf{普朗克分布与玻色-爱因斯坦统计}

式 \eqref{eq:planck} 中的 $n(\omega)$ 本质上是\textbf{玻色-爱因斯坦分布}在化学势 $\mu=0$ 时的特例：
\begin{equation}
n_\mathrm{BE}(\omega)=\frac{1}{e^{(\hbar\omega-\mu)/k_BT}-1},\qquad \text{声子：}\mu=0.
\end{equation}

\textbf{为什么声子的 $\mu=0$？} 因为声子数不守恒——晶格振动可以激发任意多个声子，也可以全部湮灭回基态。系统不固定声子总数，热平衡时化学势自然为零。这与光子（黑体辐射）完全相同，因此声子和光子服从同一分布，统称\textbf{普朗克分布}。

\textbf{与费米子的对比}：电子、质子等费米子服从费米-狄拉克分布
\begin{equation}
n_\mathrm{FD}(\omega)=\frac{1}{e^{(\hbar\omega-\mu)/k_BT}+1},
\end{equation}
分母中的 $+1$ 来源于泡利不相容原理（每个量子态最多一个费米子）。玻色子则``喜欢扎堆''，可以无限多个占据同一态，因此在低温下会出现玻色-爱因斯坦凝聚。

\textbf{1. 为什么高温 $N_{\text{ph}}\propto T$ 而低温 $N_{\text{ph}}\propto T^3$？}

\begin{center}
\begin{tabular}{p{3cm}p{5cm}p{5cm}}
\toprule
 & \textbf{高温} ($T\gg\Theta_D$) & \textbf{低温} ($T\ll\Theta_D$) \\
\midrule
被激发模式数 & 全部 $3N$ 个（常数） & 仅 $\omega\lesssim k_BT/\hbar$ 的模式，数 $\propto T^3$ \\
每个模式声子数 & $n\approx k_BT/\hbar\omega\gg1$，$\propto T$ & $n\sim O(1)$ \\
总声子数 & $(\text{常数})\times T\Rightarrow\propto T$ & $T^3\times O(1)\Rightarrow\propto T^3$ \\
\bottomrule
\end{tabular}
\end{center}

推导过程中已给出详细积分（高温 $e^x-1\approx x$，低温用 $2\zeta(3)$）。此表仅作快速参考：高温时全部模式都被激发，总声子数正比于每个模式的声子数（$\propto T$）；低温时只有少量低频模式被激发（模式数 $\propto T^3$），每个模式声子数 $\sim O(1)$，故总声子数 $\propto T^3$。

\textbf{2. 与热容 $C_V$ 的对比}

德拜热容的推导与声子总数类似，但热容积分的是能量 $\hbar\omega$ 而非声子数 $1$：
\begin{equation}
U=\int_0^{\omega_D}\hbar\omega\,n(\omega)D(\omega)\,d\omega\propto T^4\;(T\ll\Theta_D),
\end{equation}
故 $C_V=\partial U/\partial T\propto T^3$（德拜 $T^3$ 律）。

对比声子总数：$N_{\text{ph}}\propto T^3$（低温），$U\propto T^4$，$C_V\propto T^3$。三者的温度依赖不同，但都源于同一模式密度 $D(\omega)\propto\omega^2$ 和普朗克分布。热容比声子总数多一个能量因子 $\hbar\omega$，导致多一个 $T$ 的幂次。

\textbf{3. 各向同性为什么还能区分纵波和横波？}

\textbf{（a）``各向同性''到底同的是什么？}

这个概念困惑非常普遍。\textbf{各向同性不是说``没有方向''，而是说``物理性质不随方向变化''}。在任何一条你选定的传播方向上，介质都一视同仁地响应——但传播方向和质点振动方向之间的\textbf{相对关系}仍然是客观存在的。

想象一根均匀橡皮筋（各向同性）：
\begin{itemize}
    \item 你沿长度方向推拉它（纵波）——橡皮筋被压缩/拉伸，长度改变，\textbf{体积变化}，感受到的阻力由体积模量 $B$ 和剪切模量 $\mu$ 共同决定
    \item 你横向抖动它（横波）——橡皮筋发生侧向弯曲，各截面横向错动，\textbf{形状改变但长度不变}（体积近似不变），感受到的阻力仅由剪切模量 $\mu$ 决定
\end{itemize}

无论你从哪个方向观察这根橡皮筋，上述两种响应方式都存在，且各自的波速固定不变——这就是各向同性。如果橡皮筋是各向异性的（比如有纤维取向），那么不同方向上的波速会不同，但纵/横两种偏振模式依然存在。

\textbf{（b）为什么横波是纯剪切？——从形变类型理解}

固体中的弹性形变分为两种基本类型：

\begin{center}
\begin{tabular}{p{3cm}p{5cm}p{5cm}}
\toprule
 & \textbf{压缩/膨胀}（纵波） & \textbf{剪切}（横波） \\
\midrule
形变图像 & 弹簧被压缩/拉伸 & 一叠纸被横向推动，各层错开 \\
体积变化 & \textbf{改变}（$\Delta V\neq0$） & \textbf{不变}（$\Delta V=0$） \\
形状变化 & 可能不变 & 一定改变 \\
恢复力来源 & 体积模量 $B$ + 剪切模量 $\mu$ & 仅剪切模量 $\mu$ \\
\bottomrule
\end{tabular}
\end{center}

横波传播时，介质质点在垂直于传播方向的方向上振动。设想一列横波沿 $x$ 方向传播：
\begin{itemize}
    \item $x$ 处质点向上位移，$x+dx$ 处质点向下位移
    \item 两相邻层面之间发生横向错动，像一叠扑克牌被横向推开
    \item 这种错动不改变层与层之间的距离（$dx$ 不变），故体积不变
    \item 但层间相对位移产生了剪切应力，提供恢复力
\end{itemize}

这就是``纯剪切''的物理：只改变形状，不改变体积。

\textbf{（c）Helmholtz 分解的物理}

数学上，任意位移场 $\vec u(\vec r,t)$ 可分解为无旋部分和无散部分：
\begin{equation}
\vec u=\vec u_l+\vec u_t,\qquad \nabla\times\vec u_l=0,\quad \nabla\cdot\vec u_t=0.
\end{equation}

这种分解之所以可行，是因为压缩形变和剪切形变在能量上是\textbf{独立}的——改变体积不做剪切功，改变形状不做体积功。两种形变对应两种独立的弹性恢复力，故产生两种独立的波：
\begin{equation}
v_l=\sqrt{\frac{B+\frac{4}{3}\mu}{\rho}},\qquad v_t=\sqrt{\frac{\mu}{\rho}}.
\end{equation}

横波有两个正交偏振方向（均垂直于 $\vec q$），各向同性下二者简并（波速相同），故共\textbf{1支纵波 + 2支横波}。

\textit{与真实晶体的对比}：在各向异性晶体中，纵波和横波的波速还会随传播方向变化，横波的两个偏振也可能不简并。德拜模型的``各向同性''假设把方向依赖性平均掉了，但保留了偏振（纵/横）这个根本区别——因为压缩和剪切是两种物理上完全不同的形变模式。

\textbf{4. 平均声速 $v_p$ 的定义}

注意 $v_p$ \textbf{不是}简单的算术平均或调和平均，而是
\begin{equation}
\frac{3}{v_p^3}=\frac{1}{v_l^3}+\frac{2}{v_t^3}.
\end{equation}

这是由模式密度对 $v^{-3}$ 加权决定的。由于低频声学支中纵波和横波各贡献一定比例的模式，$v_p$ 实际上是\textbf{对模式密度的倒立方平均}。不要误写成 $3/v_p=1/v_l+2/v_t$。

\textbf{5. 德拜温度的物理意义}

$\Theta_D$ 是德拜模型的特征温度，标志着：
\begin{itemize}
    \item $T\gg\Theta_D$：经典极限，能量均分定理适用，$C_V\to3Nk_B$
    \item $T\ll\Theta_D$：量子极限，只有少量低频声子被激发，$C_V\propto T^3$
    \item 同一数量级：$\Theta_D$ 越高，晶格越``硬''（原子结合强、声速大），热激发越难
\end{itemize}

典型金属的德拜温度：Cu $\sim340\,\mathrm{K}$，Al $\sim428\,\mathrm{K}$，金刚石 $\sim2230\,\mathrm{K}$。金刚石的高 $\Theta_D$ 解释了它为什么热导率极高、比热极低（室温下 $T\ll\Theta_D$）。

\textbf{6. 面心立方的 $N/V$}

面心立方晶胞边长为 $a$，顶点 8 个原子各贡献 $1/8$，面心 6 个原子各贡献 $1/2$：
\begin{equation}
N_{\text{cell}}=8\times\frac{1}{8}+6\times\frac{1}{2}=1+3=4.
\end{equation}

这是 FCC 结构的基础知识，考场上必须熟记，不要临时数原子。
\end{insightbox}

\begin{thinkbox}[考场速记：德拜模型问题的五步法]
遇到``德拜模型求声子数/热容/德拜温度''类问题：
\begin{enumerate}
    \item \textbf{写模式密度}：$D(\omega)=\dfrac{3V\omega^2}{2\pi^2v_p^3}$，记住 $3/v_p^3=1/v_l^3+2/v_t^3$
    \item \textbf{写声子数/能量积分}：$N_{\text{ph}}=\int nD\,d\omega$ 或 $U=\int\hbar\omega\,nD\,d\omega$
    \item \textbf{无量纲化}：令 $x=\hbar\omega/k_BT$，积分化为 $\int x^2/(e^x-1)\,dx$ 或其变体
    \item \textbf{取极限}：高温 $x\ll1$ 展开 $e^x-1\approx x$；低温 $x_D\to\infty$ 用定积分 $2\zeta(3)$ 或 $\pi^4/15$
    \item \textbf{定常数}：德拜温度由 $3N=\int_0^{\omega_D}D(\omega)\,d\omega$ 确定
\end{enumerate}
\textit{记忆口诀}：模式密度立方二，普朗克分布积分替，高温展开低温限，德拜温度截断定。
\end{thinkbox}
